\section{True eventual orders on the same original surface}

Keep the original Noetherian local domain $S$, the original parameters
$\mathfrak m=(\pi,z)$, the original prime property of $\pi$, and
$z\notin(\pi)$. For every $N>0$ use the already constructed SAME actual DVR
$V_N=(S[U]/(\pi U-z^N))_{\mathfrak mB_N}$ and its actual original coefficient
map $j_N$. No valuation family, DVR, factorization, leading coefficient,
desired order, or slope formula is supplied as an input.

\begin{definition}[Actual original residual and vertical rings]
\label{def:gauss-surface-original-residual-vertical-rings}
\uses{def:gauss-regular-surface-weighted-chart}
\lean{EnvelopeProperty.GaussRegularSurfaceValuationOrders.residualRing,
EnvelopeProperty.GaussRegularSurfaceValuationOrders.residualCoefficientMap,
EnvelopeProperty.GaussRegularSurfaceValuationOrders.residualParameter,
EnvelopeProperty.GaussRegularSurfaceValuationOrders.verticalRing,
EnvelopeProperty.GaussRegularSurfaceValuationOrders.verticalCoefficientMap,
EnvelopeProperty.GaussRegularSurfaceValuationOrders.verticalParameter}
\leanok
Put $\overline S=S/(\pi)$ and $T=S_{(\pi)}$, with their actual original
quotient and localization coefficient maps. Write $\bar z$ and $\pi_T$ for
the actual images of the original $z$ and $\pi$. The ideal $(\pi)$ is prime
by the original prime-element hypothesis, so these rings have their derived
domain structures. They retain every original coefficient.
\end{definition}
\begin{proof}
Use exactly the original principal ideal $(\pi)$, the actual ideal quotient,
and the actual localization at its prime complement. The ideal is prime
because $\pi$ is a nonzero prime element of the original domain. Quotients
and localizations of the original Noetherian ring inherit Noetherianity.
The original quotient is nontrivial and is local by the surjective quotient
map from the original local ring. No substitute ring or chosen residue-field
isomorphism is used.
\end{proof}

\begin{lemma}[The two internally derived original DVRs]
\label{lem:gauss-surface-original-residual-vertical-dvrs}
\uses{def:gauss-surface-original-residual-vertical-rings}
\lean{EnvelopeProperty.GaussRegularSurfaceValuationOrders.residual_maximalIdeal,
EnvelopeProperty.GaussRegularSurfaceValuationOrders.residual_isDVR,
EnvelopeProperty.GaussRegularSurfaceValuationOrders.residualParameter_irreducible,
EnvelopeProperty.GaussRegularSurfaceValuationOrders.vertical_maximalIdeal,
EnvelopeProperty.GaussRegularSurfaceValuationOrders.vertical_isDVR,
EnvelopeProperty.GaussRegularSurfaceValuationOrders.verticalParameter_irreducible}
\leanok
The SAME $\overline S$ and $T$ are actual DVRs. Their genuine maximal ideals
are respectively $(\bar z)$ and $(\pi_T)$, so the stated original images
are true uniformizers for their native additive valuations.
\end{lemma}
\begin{proof}
The quotient map takes the original maximal ideal to the quotient maximal
ideal. Apply it to the original equality $\mathfrak m=(\pi,z)$; its actual
image of $\pi$ is zero, so this maximal ideal is precisely $(\bar z)$.
The original condition $z\notin(\pi)$ says that $\bar z\ne0$. Thus the
quotient local domain is not a field. Its Noetherianity and principal maximal
ideal derive that this actual quotient is a DVR, with $\bar z$ a uniformizer.

The actual maximal ideal in $T$ is the localization of the SAME $(\pi)$,
hence is precisely $(\pi_T)$. The original nonzero $\pi$ stays nonzero by
localization injectivity. Since $\pi_T$ belongs to that actual maximal ideal,
this local domain is not a field either. The same Noetherian principal
maximal-ideal criterion derives its DVR structure and the uniformizer claim.
\end{proof}

\begin{lemma}[Internally obtained original leading expansion]
\label{lem:gauss-surface-original-leading-expansion}
\uses{lem:gauss-surface-original-residual-vertical-dvrs}
\lean{EnvelopeProperty.GaussRegularSurfaceValuationOrders.exists_leadingExpansion}
\leanok
\leavevmode\par
For EVERY nonzero original $f\in S$, there are internally produced
$a,b\in\mathbf N$ and original $f',u,g\in S$ satisfying
\[
 f=\pi^a f',\qquad \pi\nmid f',\qquad
 f'=z^b u+\pi g,\qquad u\in S^\times.
\]
Their values are actual native orders:
$\operatorname{ord}_T(f)=a$ and
$\operatorname{ord}_{\overline S}(\overline{f'})=b$.
No factorization or leading-order certificate is an input.
\end{lemma}
\begin{proof}
Noetherianity of the ORIGINAL domain gives the ascending-chain condition on
principal ideals, hence finite multiplicity of the original nonunit prime
$\pi$ in the nonzero $f$. Take its actual finite maximal power to obtain
$f=\pi^af'$ with $\pi\nmid f'$. This argument needs no UFD hypothesis.

The true image $\overline{f'}$ is nonzero in the actual quotient DVR.
Its native DVR factorization gives
$\overline{f'}=\bar z^b\bar u$ for a genuine quotient unit $\bar u$.
Lift this actual element through the original surjective quotient map.
The map is a local homomorphism, so the original lift $u$ is itself a unit.
The actual kernel is exactly $(\pi)$, and therefore
$f'-z^bu=\pi g$ for an original $g\in S$.

In the original vertical localization, $f'$ is a unit because it avoids
the original ideal $(\pi)$. Its native order is zero, while the order of
the actual $\pi_T$ is one; applying the true additive valuation to the
original factorization gives $\operatorname{ord}_T(f)=a$.
In the actual quotient, the displayed unit times $\bar z^b$ similarly has
native order $b$. These are identities for the original rings and maps.
\end{proof}

\begin{theorem}[Actual eventual native orders for all original functions]
\label{thm:gauss-surface-original-eventual-orders}
\uses{lem:gauss-surface-original-leading-expansion,thm:gauss-regular-surface-original-dvr}
\lean{EnvelopeProperty.GaussRegularSurfaceValuationOrders.exists_eventual_addVal}
\leanok
For EVERY nonzero original $f\in S$, internally choose the same $a,b$ as
above. For every original $N>b$ with $N>0$, the native additive valuation of
the SAME constructed original DVR satisfies
\[
 \operatorname{ord}_{V_N}(j_N(f))=aN+b,
 \qquad \operatorname{ord}_T(f)=a.
\]
The threshold and both coefficients are derived, rather than supplied as
an order comparison or valuation-family certificate.
\end{theorem}
\begin{proof}
Obtain the actual original leading expansion just proved. Since $u$ is an
original unit, its true image under every original coefficient map $j_N$ is
a unit and has native order zero. The already constructed $V_N$ has native
orders $\operatorname{ord}_{V_N}(z)=1$ and
$\operatorname{ord}_{V_N}(\pi)=N$. Thus the actual first summand $z^bu$ has
order $b$. The actual remainder $\pi g$ has order at least $N$, since the
native order of the image of every element of $S$ is nonnegative.
When $N>b$ these orders are strictly distinct. The genuine additive
valuation equality for a sum of distinct values shows that
$\operatorname{ord}_{V_N}(f')=b$. Apply multiplicativity to
the ORIGINAL $f=\pi^af'$ to obtain $aN+b$. The vertical-order identity is
the one internally proved for this same expansion, not a new identification.
\end{proof}

This first eventual-order checkpoint concerns actual original elements of
$S$. Passing to all fractions in the SAME field, identifying the native
leading-residue description under a canonical vertical-residue field map,
normalizing by the actual $\operatorname{ord}_{V_N}(t)=eN+c$, obtaining fixed
Gauss-valued points and their true convergence, and residual difference
quotients are further statements. In particular, no assumption that
$t$ is a unit times a pure power at every original centre is made here.
